\documentclass[dvips, a4paper, listsleft, 12pt, openany, oneside]{scrbook}
% new koma package
\usepackage{scrpage2}

\usepackage{ae}
\usepackage[latin1]{inputenc}
\usepackage[ngerman]{babel}
\usepackage[T1]{fontenc}
\usepackage{verbatim}
% \usepackage[small]{caption2}
\usepackage{graphicx,psfrag}
\usepackage{amsmath}
\usepackage{amsfonts}
\usepackage{amssymb}
\usepackage[dvips]{color}
\usepackage{xspace}
\usepackage{bm}
% \usepackage{bbm}
% \usepackage{slashbox}
\usepackage{booktabs} % extends the tables property by: \toprule, \midrule and \bottomrule
\usepackage{longtable}
\usepackage{array}


%%%%%%%%%%%%%%%%%%%%%%%%%%%%%%%%%%%%%%%%%%%%%%%%%%%%%%%%%%%%
% if hyperrefs are desired, uncomment the appropriate lines in the following block;

% define your own link color(including: \cite, \ref, \eqref, \footnote, ...)
% \definecolor{darkblue}{rgb}{0,0,.62} 
% or:
% \definecolor{darkblue}{rgb}{0,.19,.40}

% \usepackage[breaklinks, colorlinks, urlcolor=darkblue]{hyperref}
% \usepackage{hypcap}
% NOTE: sometimes problems arise when using hyperref;
%       especially the dvi/ps file might look different from what you expected;
\usepackage{url}

%%%%%%%%%%%%%%%%%%%%%%%%%%%%%%%%%%%%%%%%%%%%%%%%%%%%%%%%%%%%
% for koma headers
\pagestyle{scrheadings}
% define head and foot marks
\renewcommand{\chaptermark}[1]{\markboth{#1}{}}
\renewcommand{\sectionmark}[1]{\markright{-~#1}{}}
\ihead{\leftmark}
\chead{}
\ohead{\pagemark}
\cfoot{}
\setheadsepline{.4pt}

%%%%%%%%%%%%%%%%%%%%%%%%%%%%%%%%%%%%%%%%%%%%%%%%%%%%%%%%%%%%
% chapter head change
% adding lines above the chapter head
% 1st get a new command
\newcommand*{\ORIGchapterheadstartvskip}{}%
% 2nd save the original definition to the new command
\let\ORIGchapterheadstartvskip=\chapterheadstartvskip
% 3rd redefine the command using the saved original command
\renewcommand*{\chapterheadstartvskip}{%
  \ORIGchapterheadstartvskip
  {%
    \setlength{\parskip}{0pt}%
    \noindent\rule[0.3\baselineskip]{\linewidth}{0.8pt}\par
  }%
}
% and below the chapter head
\newcommand*{\ORIGchapterheadendvskip}{}%
\let\ORIGchapterheadendvskip=\chapterheadendvskip
\renewcommand*{\chapterheadendvskip}{%
  {%
    \setlength{\parskip}{0pt}%
    \noindent\rule[0.3\baselineskip]{\linewidth}{0.8pt}\par
  }%
  \ORIGchapterheadendvskip
}

% %%%%%%%%%%%%%%%%%%%%%%%%%%%%%%%%%%%%%%%%%%%%%%%%%%%%%%%%%%%%
% % list of abbreviations
% \usepackage{nomencl}
% % rename command to 'abk'
% \let\abk\nomenclature
% % title in german
% \renewcommand{\nomname}{Abk�rzungsverzeichnis}
% % define space between abbreviation and declaration
% \setlength{\nomlabelwidth}{.25\hsize}
% \renewcommand{\nomlabel}[1]{#1 \dotfill}
% % decrease row spacing
% \setlength{\nomitemsep}{-\parsep}
% \makeglossary
% % new command: highlighting abbrevation letters in declaration part
% \usepackage[normalem]{ulem}
% \newcommand{\markup}[1]{\textbf{#1}}
% 
% %%%%%%%%%%%%%%%%%%%%%%%%%%%%%%%%%%%%%%%%%%%%%%%%%%%%%%%%%%%%
% % define german names
%   \def\prefacename{Vorwort}%
%   \def\refname{Literatur}%
%   \def\abstractname{Zusammenfassung}%
%   \def\bibname{Literaturverzeichnis}%
%   \def\chaptername{Kapitel}%
%   \def\appendixname{Anhang}%
%   \def\contentsname{Inhaltsverzeichnis}% % oder nur: Inhalt
%   \def\listfigurename{Abbildungsverzeichnis}%
%   \def\listtablename{Tabellenverzeichnis}%
%   \def\indexname{Index}%
%   \def\figurename{Abbildung}%
%   \def\tablename{Tabelle}%  % or: Tafel
%   \def\partname{Teil}%
%   \def\enclname{Anlage(n)}% % or: Beilage(n)
%   \def\ccname{Verteiler}%   % or: Kopien an
%   \def\headtoname{An}%
%   \def\pagename{Seite}%
%   \def\seename{siehe}%
%   \def\alsoname{siehe auch}

%%%%%%%%%%%%%%%%%%%%%%%%%%%%%%%%%%%%%%%%%%%%%%%%%%%%%%%%%%%%
% define pagesize
\setlength{\oddsidemargin}{5mm}
\setlength{\evensidemargin}{5mm}
\setlength{\textwidth}{155mm}
\setlength{\topmargin}{-5mm}
\setlength{\textheight}{235mm}

%%%%%%%%%%%%%%%%%%%%%%%%%%%%%%%%%%%%%%%%%%%%%%%%%%%%%%%%%%%%
% special formatted strings
\newcommand{\matlab}{\emph{MATLAB}\xspace}
\newcommand{\intel}{\emph{INTEL}\xspace}

%%%%%%%%%%%%%%%%%%%%%%%%%%%%%%%%%%%%%%%%%%%%%%%%%%%%%%%%%%%%
% changing some fonts 
% \renewcommand{\captionlabelfont}{\small\sffamily\bfseries}
% \setcaphanging

\setkomafont{pagehead}{%
\normalfont\sffamily\bfseries
}
\setkomafont{pagenumber}{%
\normalfont\rmfamily\slshape
}

%%%%%%%%%%%%%%%%%%%%%%%%%%%%%%%%%%%%%%%%%%%%%%%%%%%%%%%%%%%%
% command for the list of symbols
\newcommand{\listsymbol}[2]{#1$\;\;$\dotfill$\;\;$&#2\\}

% define graficspath

\graphicspath{{bilder/}{bilder/Kapitel_1/}{bilder/Kapitel_2/}}
\begin{document}
%%%%%%%%%%%%%%%%%%%%
% the customized title page
% this is the title page 
\begin{titlepage}
\begin{center}
  \begin{figure}
    \begin{center}
      \includegraphics[scale=0.5]{bilder/uni_logo.eps}
    \end{center}
  \end{figure}
\normalsize{\textbf{Institut f�r Elektrotechnik und Informationstechnik}\\
Universit�t Paderborn\\
Fachgebiet Nachrichtentechnik\\
Prof. Dr.-Ing. Reinhold H�b-Umbach\\}
\vfill {\bf {\Large Seminararbeit} \linebreak \\ \LARGE{Training des statistischen Sprachmodells}}
\renewcommand{\baselinestretch}{2}\normalsize
\vfill{von\\ {\large Buyu Xiao\\ Matr.-Nr.: 6436310}}
\renewcommand{\baselinestretch}{1,5}\normalsize
\vfill{
  \begin{flushleft}
    Betreuer: Volker Leutnant\\
  \end{flushleft}
}
\renewcommand{\baselinestretch}{1}\normalsize
\end{center}
\end{titlepage}
% 
% %%%%%%%%%%%%%%%%%%%%
% % frontmatter
% \frontmatter
% \tableofcontents
% \clearpage
% % list of figures
% \addcontentsline{toc}{chapter}{Abbildungsverzeichnis}
% \listoffigures
% \clearpage
% % list of tables
% \addcontentsline{toc}{chapter}{Tabellenverzeichnis}
% \listoftables
% \clearpage
% %list of equation symbols 
% \include{symbols}
% \clearpage
% % list of abbreviations
% \addcontentsline{toc}{chapter}{Abk�rzungsverzeichnis}
% % \printglossary
% \clearpage
% \newpage
% \mbox{} \thispagestyle{empty} \newpage

%%%%%%%%%%%%%%%%%%%%
% mainmatter
\mainmatter
% make the header to include the chapter numbering
% from now on
\ihead{\thechapter~\leftmark} %~\rightmark}
% the chapters
\chapter{Kurze Einf�hrung}
\label{chapter:introduction}

In diesem Kapitel werden die wichtigsten LaTeX-Konstrukte kurz vorgestellt.\\
Das PDF-Dokument PDF \textit{short-math-guide.pdf} in dem ZIP-Ordner stellt eine gute Anlaufstelle bei Fragen und Problemen (siehe Ordner: \textsc{lesenswertes}).

\section{Formeln}
\label{section:equations}

Dieser Abschnitt befa�t sich mit der Verwendung von Formeln.

\subsection{Inline-Formeln}
\label{subsection:inlineEquation}

Inline-Formeln beginnen und enden mit dem \$-Zeichen: $e^x = \sum\limits_{n=0}^{\infty}\frac{x^n}{n!}$

\subsection{Die equation-Umgebung}
\label{subsection:defaultEquation}

Die im vorherigen Abschnitt \ref{subsection:inlineEquation} eingef�hrte Formel schreibt sich mit Hilfe der \textit{equation}-Umgebung wie folgt:
\begin{equation}
\label{equation:expFunktion}
 e^x = \sum\limits_{n=0}^{\infty}\frac{x^n}{n!}
\end{equation}
Die Exponential-Funktion aus Gleichung \eqref{equation:expFunktion} wird in der Mustererkennung h�ufig verwendet.
ACHTUNG: \textit{eqref} fa�t die Nummer der Gleichung automatisch in runde Klammern.

\section{Literaturquellen referenzieren}
\label{section:literaturref}

Die Literaturquellen werden mit Hilfe von \textsc{bibtex} referenziert. Die Quellen k�nnen in unterschiedlichen  ''.bib'' Dateien beschrieben sein (siehe z.B. \url{bibs/RefA}). Der Aufruf im Dokument erfolgt dann gem��:\\
\textbackslash bibliographystyle\{IEEEtran\}\\
\textbackslash bibliography\{IEEEfull,bibs/RefA, bibs/RefB\}\\
Eine Auflistung aller bibliographystyles findet man z.B. unter \url{http://www.cs.stir.ac.uk/~kjt/software/latex/showbst.html}. Hier wurde der \textit{IEEEtrans} style verwendet, der in der \url{bibs/IEEEtran.bst} definiert wurde. Die Abk�rzungen stehen in \url{bibs/IEEEfull.bib}, welche vor allen Bibliotheken eingebunden wurde. Alle referenzierten  Quellen werden automatisch in das Literaturverzeichnis eingf�gt, ein Fehler erscheint, wenn der \\cite\{\}-Befehl nicht verwendet wurde.\\
Eine Quelle aus einer Datenbank wird so referenziert: \textbackslash cite\{bibtexkey\}.\\
So ist \cite{Digalakis1993} z.B. aus Datenbank RefA und \cite{Lee2007} z.B. aus Datenbank RefB.

\section{Tabellen}
\label{section:tabular}

Dieser Abschnitt befa�t sich mit der Verwendung von Tabellen und den neuen Befehlen des \textit{booktabs}-Packets.
\begin{table}[h]
  \begin{center}
    \begin{tabular}{clr}
      \toprule
      \bf laufende & \bf Wert 1 & \bf Wert 2\\
      \bf Nummer  & [Einheit]  & [Einheit] \\
      \midrule
      $1$ &  $1\,000$ & $20$  \\
      $2$ &  $1\,500$ & $40$  \\
      \bottomrule
    \end{tabular}
  \end{center}
\caption{Eine einfache Tabelle}
\label{tab:table_1}
\end{table}

\section{Grafiken}
\label{section:figures}

Dieser Abschnitt stellt kurz das Einf�gen von Grafiken und das Ersetzen von Text innerhalb einer eps-Datei durch LaTeX-eigene Schriften und insbesonder auch Formeln vor. 

Die Verwendung des \textit{lesenswertes}-Packets wird in der \textit{psfrag-guide.pdf} Datei n�her erkl�rt (siehe Ordner: \textsc{usefullreadings}).\\
\\
ACHTUNG: Die Ersetzungen sind erst in der PS/PDF-Datei an der gew�nschten Stelle. In der DVI-Datei sind die Ersetzungen nur aufgelistet.
\begin{figure}[h]
  \begin{center}
    \resizebox{\textwidth}{!}
    {
    \psfrag{Hintergrund}[][c]{Hintergrund}
    \psfrag{Gesicht}[][c]{Gesicht}
    \psfrag{Detektionskaskade}[][ct]{Detektionskaskade}
    \psfrag{Analyse-}[][ct]{Analyse-}
    \psfrag{fenster}[][c]{fenster}
    \psfrag{W}[][cb]{$\mathcal{W}$}
    \psfrag{Stufe1:}[][ct]{Stufe 1:}
    \psfrag{Stufe 2:}[][ct]{Stufe 2:}
    \psfrag{Stufe 3:}[][ct]{Stufe 3:}
    \psfrag{Stufe 4:}[][ct]{Stufe 4:}
    \psfrag{20 Pixel}[][c]{$|\mathcal{W}_1'|$=$20$ Pixel} 
    \psfrag{40 Pixel}[][c]{$|\mathcal{W}_2'|$=$40$ Pixel} 
    \psfrag{160 Pixel}[][c]{$|\mathcal{W}_3'|$=$160$ Pixel} 
    \psfrag{244 Pixel}[][c]{$|\mathcal{W}_4'|$=$244$ Pixel} 
    \psfrag{T1=}[][c]{$T_1=$}
    \psfrag{T2=}[][c]{$T_2=$}
    \psfrag{T3=}[][c]{$T_3=$}
    \psfrag{0.345}[][c]{$0.345$}
    \psfrag{0.375}[][c]{$0.375$}
    \psfrag{0.385}[][c]{$0.385$}
    \psfrag{T4}[][c]{$T_4$}
    \includegraphics{beispiel.eps}
    }
  \end{center}
  \caption{Eine Grafik mit ersetzten Schriften}
  \label{fig:figure_1}
\end{figure} 





% the chapters
\chapter{Grafiken die 2.}
\label{chapter:grafiken2}

Der Ort, an dem die Grafiken sich befinden, kann bei der Verwendung des \textit{graphicx} Pakets mit Hilfe des Befehls 
\begin{verbatim}
\graphicspath{{bilder/} {bilder/chapter1/} {bilder/chapter2/}}} 
\end{verbatim}
dem LaTeX-Interpreter vor Begin des Dokuments aber auch an jeder beliebigen Stelle im Dokument neu mitgeteilt werden. Dabei ist darauf zu achten, dass die erneute Verwendung des Befehls die vorherigen Werte �berschreibt. Der Pafd kann alternativ dazu aber auch direkt (relativ oder absolut) dem Befehl \slash includegraphics mitgeteilt werden.
\begin{figure}[h]
  \begin{center}
    \resizebox{0.6\textwidth}{!}{
    \includegraphics{word_phones_states_mixdens.eps}
    }
  \end{center}
  \caption{Verwendung des Pfades der durch \slash graphicspath gegeben ist}
  \label{fig:figure_2}
\end{figure} 
\begin{figure}[h]
  \begin{center}
    \resizebox{0.6\textwidth}{!}{
    \includegraphics{bilder/Kapitel_2/word_phones_states_mixdens.eps}
    }
  \end{center}
  \caption{Verwendung des relativen Pfades}
  \label{fig:figure_3}
\end{figure} 






% ... more to come, hopefully!
\clearpage

%%%%%%%%%%%%%%%%%%%%
% this is where the backmatter starts
% \backmatter % toggling this comment results in chapter names with or without capital numbering
% the bibliography
\addcontentsline{toc}{chapter}{Literaturverzeichnis}
% the bibliography shouldn't have a consecutive numbering
\ihead{\leftmark}
\bibliographystyle{bibs/IEEEtran}
\bibliography{bibs/IEEEfull,bibs/RefA,bibs/RefB}
\clearpage
% but the appendix could
\ihead{\thechapter~\leftmark}
% the appendix is part of the backmatter, but uses capital numbering by default
\appendix
\begin{appendix}
\chapter{Inhalts�bersicht der Begleit-DVD}
Alle Zusatzinformationen sind in digitaler Form auf der beiliegenden DVD enthalten. Im Folgenden wird die Verzeichnisstruktur der DVD erl�utert:
\begin{itemize}
\setlength{\itemsep}{0pt}
  \item \textbf{Ordner 1} -- Enth�lt ...
  \begin{itemize}
    \item[$\bullet$] \textbf{Unterordner 1.1} -- Enth�lt ... .
    \item[$\bullet$] \textbf{Unterordner 2.1} -- Enth�lt ... .
    \item[$\bullet$] \textbf{Unterordner 3.1} -- Enth�lt ... .
    \item[$\bullet$] \textbf{Unterordner 4.1} -- Enth�lt ... .
  \end{itemize}
  \item \textbf{Ordner 2} -- ... .
  \item \textbf{Ordner 3} -- ... .
  \item \textbf{Ordner 4} -- ... .
  \item \textbf{Ordner 5} -- ... .
  \item \textbf{Ordner 6} -- Enth�lt ... :
  \begin{itemize}
    \item[$\bullet$] \textbf{Unterordner 6.1} -- Enth�lt ... .
    \item[$\bullet$] \textbf{Unterordner 6.2} -- Enth�lt ... .
  \end{itemize}
  \item \textbf{Ordner 7} -- Enth�lt ... .
\end{itemize}
Der Inhalt der aufgef�hrten Verzeichnisse ist zur besseren �bersicht ggf. in Unterverzeichnisse aufgeteilt.
\end{appendix}


\end{document}


%%%%%%%%%%%%%%%%%%%%
% this is how to configure kile to take advantage
% of the nomenclature package

% configure kile --> quick build:
% LATEX
% BIBTEX
% LATEX
% LATEX
% VIEWDVI

% use iso 8859-1 or iso 8859-15 to write your text!